\documentclass[12pt]{article}

\usepackage[a4paper,margin=1in]{geometry}
\usepackage{amsmath,amssymb}
\usepackage{newtxtext,newtxmath}
\usepackage{fancyhdr}
\usepackage{listings}
\usepackage{graphicx}
\usepackage{float}
\usepackage{booktabs}
\usepackage{xcolor}
\usepackage{hyperref}

\setlength{\parindent}{0pt}
\setlength{\parskip}{5pt}

\AtBeginDocument{
    \fontsize{12pt}{14pt}\selectfont
}

% Header
\pagestyle{fancy}
\fancyhf{}
\fancyhead[L]{Name: \underline{Shaik Jnaved}}
\fancyhead[R]{Enrollment No.: \underline{P26DS001}}
\renewcommand{\headrulewidth}{0.4pt}
\setlength{\headheight}{15pt}

% Code formatting
\lstset{
    language=C++,
    basicstyle=\ttfamily\fontsize{8pt}{10pt}\selectfont,
    breaklines=true,
    frame=single,
    columns=fullflexible,
    keepspaces=true,
    showstringspaces=false,
    keywordstyle=\color{blue}\textbf,
    commentstyle=\color{green!50!black}\textit,
    tabsize=4,
    numbers=left,
    numberstyle=\tiny\color{gray},
    numbersep=6pt
}

\begin{document}

% ==============================================================================
% TITLE & HEADER
% ==============================================================================
\begin{center}
\textbf{Computer Science and Engineering Department, SVNIT, Surat}\\
\textbf{M.Tech. I -- Semester I}\\
\textbf{High Performance Computing}\\
\textbf{Lab Assignment 7}
\end{center}

\vspace{0.2cm}

% ==============================================================================
% PROBLEM 1: TASK ALLOCATION PROBLEM (CUDA C vs OpenMP vs Sequential)
% ==============================================================================
\textbf{\large Problem Statement 1}

\textbf{Title:} Parallel Metaheuristics (GA \& PSO) for Task Allocation Problem using CUDA C, OpenMP, and Sequential C.

\textbf{Objective:}
Implement CUDA C parallel versions of Genetic Algorithm (GA) and Particle Swarm Optimization (PSO) to solve the Task Allocation Problem (TAP) discussed in Lab Assignment 5 ($M = 300$ tasks, $P = 10$ processing nodes). Minimize Total Cost comprising Execution Cost ($\sum E[i][a_i]$) and Communication Cost ($\sum_{i < j} C[i][j] \cdot [a_i \ne a_j]$). Compare the running time performance of CUDA C against the sequential and OpenMP implementations on an NVIDIA RTX 4050 GPU.

\vspace{0.2cm}

\textbf{Source Code (\texttt{lab\_7/tap\_solver.cu}):}
\lstinputlisting{../../lab_7/tap_solver.cu}

\newpage

\textbf{Performance Analysis \& Discussion (Task Allocation Problem):}

\begin{figure}[H]
\centering
\includegraphics[width=0.95\textwidth]{images/tap_performance_comparison.png}
\caption{Task Allocation Problem: (a) Execution Time Comparison and (b) Parallel Speedup on NVIDIA RTX 4050 GPU vs. 16-Thread OpenMP vs. Sequential CPU.}
\end{figure}

\begin{table}[H]
\centering
\caption{Task Allocation Problem Performance Summary ($M=300, P=10, N=100, 100\text{ Iterations}$)}
\vspace{0.2cm}
\begin{tabular}{lcccc}
\toprule
\textbf{Algorithm \& Version} & \textbf{Time (s)} & \textbf{Best Cost} & \textbf{Speedup vs. Seq} & \textbf{Platform} \\
\midrule
GA -- Sequential C           & 0.3306 s & 198,210.00 & 1.00$\times$  & Host CPU (1 core) \\
GA -- OpenMP Parallel        & 0.0697 s & 209,539.00 & 4.75$\times$  & Host CPU (16 threads) \\
GA -- CUDA C Parallel        & \textbf{0.0247 s} & 198,769.00 & \textbf{13.39$\times$} & RTX 4050 GPU \\
\midrule
PSO -- Sequential C          & 0.3632 s & 206,752.00 & 1.00$\times$  & Host CPU (1 core) \\
PSO -- OpenMP Parallel       & 0.0709 s & 218,886.00 & 5.12$\times$  & Host CPU (16 threads) \\
PSO -- CUDA C Parallel       & \textbf{0.0619 s} & 206,828.00 & \textbf{5.87$\times$}  & RTX 4050 GPU \\
\bottomrule
\end{tabular}
\end{table}

\textbf{Key Observations:}
\begin{enumerate}
    \item \textbf{Parallel Fitness Evaluation:} Fitness evaluation of $M=300$ tasks involves $M(M-1)/2 = 44,850$ pairwise communication checks per individual. On the GPU, a block of 128 threads per individual parallelizes the double loop and uses warp shuffle reductions (\texttt{\_\_shfl\_down\_sync}), reducing evaluation time to under $0.06$ ms per generation.
    \item \textbf{Speedup Results:} CUDA GA achieved a \textbf{13.39$\times$ speedup} over sequential C and a \textbf{2.82$\times$ speedup} over 16-thread OpenMP. CUDA PSO achieved a \textbf{5.87$\times$ speedup} over sequential C.
    \item \textbf{Solution Quality:} Both GPU implementations found competitive minimum costs ($\approx 198,769$ for GA and $\approx 206,828$ for PSO) closely matching sequential and OpenMP baselines.
\end{enumerate}

\newpage

% ==============================================================================
% PROBLEM 2: BENCHMARK FUNCTIONS (GA & PSO Across Versions)
% ==============================================================================
\textbf{\large Problem Statement 2}

\textbf{Title:} Performance Evaluation of Parallel Metaheuristics Across Five Standard Benchmark Functions.

\textbf{Objective:}
Test and benchmark the performance of the metaheuristics (Genetic Algorithm and Particle Swarm Optimization) across five classical continuous optimization benchmark functions ($D = 30$, Population = 1024, 500 Iterations) on Sequential C, OpenMP, and CUDA C architectures:
\begin{enumerate}
    \item \textbf{Sphere Function ($f_1$):} $f_1(\mathbf{x}) = \sum_{i=0}^{D-1} x_i^2$, $x_i \in [-5.12, 5.12]$, $f^* = 0.0$ (Unimodal, Convex).
    \item \textbf{Rosenbrock Function ($f_2$):} $f_2(\mathbf{x}) = \sum_{i=0}^{D-2} [100(x_{i+1} - x_i^2)^2 + (x_i - 1)^2]$, $x_i \in [-5, 10]$, $f^* = 0.0$ (Valley-shaped, Non-convex).
    \item \textbf{Rastrigin Function ($f_3$):} $f_3(\mathbf{x}) = 10D + \sum_{i=0}^{D-1} [x_i^2 - 10\cos(2\pi x_i)]$, $x_i \in [-5.12, 5.12]$, $f^* = 0.0$ (Highly Multimodal).
    \item \textbf{Ackley Function ($f_4$):} $f_4(\mathbf{x}) = -20\exp\left(-0.2\sqrt{\frac{1}{D}\sum x_i^2}\right) - \exp\left(\frac{1}{D}\sum \cos(2\pi x_i)\right) + 20 + e$, $x_i \in [-32.768, 32.768]$, $f^* = 0.0$ (Multimodal Steep Funnel).
    \item \textbf{Griewank Function ($f_5$):} $f_5(\mathbf{x}) = \frac{1}{4000}\sum x_i^2 - \prod \cos\left(\frac{x_i}{\sqrt{i+1}}\right) + 1$, $x_i \in [-600, 600]$, $f^* = 0.0$ (Multimodal Coupled).
\end{enumerate}

\vspace{0.2cm}

\textbf{Source Code (\texttt{lab\_7/benchmarks.cu}):}
\lstinputlisting{../../lab_7/benchmarks.cu}

\newpage

\textbf{Performance Analysis \& Discussion (Benchmark Functions):}

\begin{figure}[H]
\centering
\includegraphics[width=0.95\textwidth]{images/benchmark_runtime_comparison.png}
\caption{Execution Runtime Comparison for PSO and GA across 5 Benchmark Functions ($D=30, N=1024, 500\text{ Iterations}$).}
\end{figure}

\begin{figure}[H]
\centering
\includegraphics[width=0.95\textwidth]{images/benchmark_speedup_comparison.png}
\caption{Measured Parallel Speedup Factors for OpenMP (16 threads) and CUDA C (RTX 4050 GPU) across 5 Benchmark Functions.}
\end{figure}

\begin{table}[H]
\centering
\caption{Comprehensive Benchmark Performance Summary ($D = 30$, Population = 1024, 500 Iterations)}
\vspace{0.2cm}
\resizebox{\textwidth}{!}{
\begin{tabular}{llccccc}
\toprule
\textbf{Function} & \textbf{Algorithm} & \textbf{Seq Time (s)} & \textbf{OpenMP Time (s)} & \textbf{CUDA Time (s)} & \textbf{OpenMP Speedup} & \textbf{CUDA Speedup} \\
\midrule
Sphere ($f_1$)      & PSO & 0.0387 s & 0.0145 s & 0.0347 s & 2.68$\times$ & 1.12$\times$ \\
                    & GA  & 0.0336 s & 0.0246 s & 0.0684 s & 1.36$\times$ & 0.49$\times$ \\
\midrule
Rosenbrock ($f_2$)  & PSO & 0.0374 s & 0.0161 s & 0.0302 s & 2.32$\times$ & 1.24$\times$ \\
                    & GA  & 0.0374 s & 0.0310 s & 0.0558 s & 1.21$\times$ & 0.67$\times$ \\
\midrule
Rastrigin ($f_3$)   & PSO & 0.1513 s & 0.0615 s & 0.0342 s & 2.46$\times$ & \textbf{4.42$\times$} \\
                    & GA  & 0.1439 s & 0.0383 s & 0.0608 s & 3.76$\times$ & \textbf{2.37$\times$} \\
\midrule
Ackley ($f_4$)      & PSO & 0.1458 s & 0.0207 s & 0.0352 s & 7.05$\times$ & \textbf{4.14$\times$} \\
                    & GA  & 0.2131 s & 0.0672 s & 0.0612 s & 3.17$\times$ & \textbf{3.48$\times$} \\
\midrule
Griewank ($f_5$)    & PSO & 0.1735 s & 0.0338 s & 0.0401 s & 5.14$\times$ & \textbf{4.32$\times$} \\
                    & GA  & 0.2822 s & 0.0519 s & 0.0668 s & 5.43$\times$ & \textbf{4.23$\times$} \\
\bottomrule
\end{tabular}
}
\end{table}

\textbf{Conclusions:}
\begin{enumerate}
    \item \textbf{Arithmetic Intensity Matters:} For simple polynomial functions (Sphere, Rosenbrock), the low arithmetic intensity means CUDA performance is bounded by memory launch overheads. However, for functions with transcendental operations ($\cos, \exp, \sqrt{}$) like Rastrigin, Ackley, and Griewank, the massive parallel ALU throughput of the GPU yields \textbf{3.5$\times$ to 4.4$\times$ speedups}.
    \item \textbf{Convergence Quality:} PSO exhibited superior convergence speed on continuous unimodal and steep-funnel functions (reaching $2.33 \times 10^{-19}$ on Sphere and $8.35 \times 10^{-10}$ on Ackley), while Real-Coded GA showed robust performance in escaping local minima on multimodal landscapes (reaching $0.613$ on Rastrigin).
    \item \textbf{Parallel Scalability:} Both OpenMP and CUDA C implementations achieve reproducible parallel scaling without race conditions through per-thread lightweight pseudorandom generators and zero-overhead memory reduction primitives.
\end{enumerate}

\end{document}
